\chapter{Testování a~nasazení}\label{chap:testing}

\begin{chapterabstract}
	Kapitola popisuje strategii testování jednotlivých částí systému, měření pokrytí kódu, průběžnou integraci a~proces nasazení aplikace do produkčního prostředí.
\end{chapterabstract}

\section{Testování backendu}

Backend je testován na třech úrovních: unit testy izolovaných komponent, testy datové vrstvy nad in-memory databází a~integrační testy celého API. Celkem obsahuje 445 testovacích metod ve~44 testovacích třídách.

\subsection{Unit testy}

Unit testy ověřují business logiku service vrstvy a~chování kontrolerů v~izolaci. Jsou implementovány pomocí JUnit~5 a~knihovny Mockito, která umožňuje nahradit závislosti (repozitáře, mappery) simulovanými objekty (\texttt{@Mock}).

Kontrolery jsou testovány prostřednictvím \texttt{MockMvc} -- Spring testovací framework, který simuluje HTTP požadavky bez spuštění skutečného serveru. Testy ověřují správné HTTP kódy, formát odpovědí a~validaci vstupních dat. Service testy ověřují business logiku včetně chybových scénářů (neexistující entita, duplicitní identifikátor, neplatné přihlašovací údaje).

\subsection{Testy datové vrstvy}

Repozitáře a~JPA Specification jsou testovány anotací \texttt{@DataJpaTest}, která spouští pouze datovou vrstvu nad in-memory databází H2 v~PostgreSQL compatibility mode. Testy ověřují správnost dotazů odvozených z~názvů metod, explicitních \texttt{@Query} anotací a~dynamických predikátů pro filtrování karet majetku.

\subsection{Integrační testy}

Integrační testy spouštějí celou Spring Boot aplikaci (\texttt{@SpringBootTest}) včetně bezpečnostní konfigurace, serializace a~databázové vrstvy. Pokrývají kompletní API flow od autentizace přes CRUD operace až po inventurní projekci. Všechny integrační testy dědí ze~společné třídy \texttt{BaseIntegrationTest}, která zajišťuje vytvoření testovacích uživatelů a~přihlášení.

\subsection{Pokrytí kódu}

Pokrytí kódu je měřeno nástrojem JaCoCo s~následujícími minimálními prahy:

\begin{itemize}
	\item pokrytí řádků: alespoň 80\,\%,
	\item pokrytí větví: alespoň 75\,\%.
\end{itemize}

Z~měření jsou vyloučeny třídy, které neobsahují testovatelnou logiku: DTO, JPA entity, konfigurační třídy a~MapStruct mappery (generovaný kód). Konfigurace prahových hodnot je znázorněna ve~výpisu~\ref{lst:jacoco}. Při nesplnění build selže.

\begin{listing}[H]
\caption{Konfigurace JaCoCo -- minimální pokrytí kódu}\label{lst:jacoco}
\begin{minted}[fontsize=\small]{xml}
<rule>
  <element>PACKAGE</element>
  <limits>
    <limit>
      <counter>LINE</counter>
      <value>COVEREDRATIO</value>
      <minimum>0.80</minimum>
    </limit>
    <limit>
      <counter>BRANCH</counter>
      <value>COVEREDRATIO</value>
      <minimum>0.75</minimum>
    </limit>
  </limits>
</rule>
\end{minted}
\end{listing}

\section{Testování frontendu}

Frontend je testován knihovnou Vitest (7 testovacích souborů, 54 testovacích metod). Testy pokrývají kritickou business logiku -- utility funkce (formátování dat, české řazení, extrakce chybových hlášek), validační schémata a~vybrané React komponenty. Komponentové testy využívají \texttt{@testing-library/react} pro renderování a~simulaci uživatelských interakcí. Správnost UI komponent je dále zajištěna statickou typovou kontrolou TypeScriptu.

Statická typová kontrola prostřednictvím kompilátoru TypeScriptu slouží jako doplňková vrstva ověření -- při každém buildu kompilátor zkontroluje typovou konzistenci celé aplikace.

\section{Testování mobilní aplikace}

Mobilní aplikace obsahuje dva typy testů:

\begin{itemize}
	\item \textbf{Unit testy} (3 soubory) -- testují síťovou komunikaci (\texttt{ApiClient\-Test} s~knihovnou MockWebServer), synchronizační logiku (\texttt{DataSync\-Repository\-Test}) a~správu konfigurace (\texttt{BaseUrl\-Storage\-Test}). Běží na JVM bez emulátoru.
	\item \textbf{Instrumentované testy} (2 soubory) -- běží na emulátoru nebo fyzickém zařízení. Testují Room DAO operace (\texttt{DaoTest}: vkládání, dotazování, mazání, upsert) a~šifrované úložiště (\texttt{TokenStorageTest}: ukládání a~čtení tokenů, migrace z~nešifrovaného úložiště).
\end{itemize}

Uživatelské rozhraní je průběžně ověřováno manuálním testováním na cílovém zařízení Zebra TC22.

\section{Souhrn testování}

Tabulka~\ref{tab:tests} shrnuje rozsah testování všech komponent systému. U~frontendu doplňuje kompilační typová kontrola TypeScriptu dynamické testy o~statické ověření typové konzistence celé aplikace.

\begin{table}[H]\centering
\renewcommand{\arraystretch}{1.3}
\begin{tabular}{l|r|l|l}
	\textbf{Komponenta} & \textbf{Testů} & \textbf{Pokrytí kódu} & \textbf{Nástroje} \tabularnewline \hline
	Backend & 445 & $\geq$\,80\,\% ř., $\geq$\,75\,\% větví & JUnit~5, Mockito \tabularnewline
	Frontend & 54 & -- & Vitest \tabularnewline
	Mobilní apl. & 40 & -- & JUnit~5 \tabularnewline
\end{tabular}
\caption{Přehled testů jednotlivých komponent}\label{tab:tests}
\end{table}

\section{Průběžná integrace}

Projekt využívá GitHub Actions pro automatizované sestavení, testování a~statickou analýzu kódu. Workflow \texttt{sonarcloud.yml} se spouští při každém push do větve \texttt{main} a~sestává ze~čtyř kroků:

\begin{enumerate}
	\item \textbf{Backend} -- sestavení a~spuštění testů příkazem \texttt{mvn clean verify}. Výstupem jsou zkompilované třídy, výsledky testů a~JaCoCo report.
	\item \textbf{Android} -- sestavení a~spuštění unit testů. Instrumentované testy vyžadující emulátor nejsou v~CI zahrnuty a~jsou spouštěny lokálně.
	\item \textbf{Frontend} -- instalace závislostí (\texttt{npm ci}) a~spuštění testů (\texttt{npm run test}).
	\item \textbf{SonarCloud} -- po úspěšném dokončení předchozích kroků se spustí statická analýza kódu nad všemi třemi moduly. SonarCloud vyhodnocuje kvalitu kódu, bezpečnostní zranitelnosti a~duplicity.
\end{enumerate}

Zkrácená ukázka CI workflow je ve~výpisu~\ref{lst:ci}.

\begin{listing}[H]
\caption{GitHub Actions workflow pro průběžnou integraci (zkráceno)}\label{lst:ci}
\begin{minted}[fontsize=\small]{yaml}
name: SonarCloud Analysis
on:
  push:
    branches: [main]
jobs:
  backend:
    runs-on: ubuntu-latest
    steps:
      - uses: actions/checkout@v4
      - uses: actions/setup-java@v4
        with: { java-version: '21' }
      - run: mvn -B -f pom.xml clean verify
  android:
    runs-on: ubuntu-latest
    steps:
      - run: ./gradlew assembleDebug testDebugUnitTest
  frontend:
    runs-on: ubuntu-latest
    steps:
      - run: npm ci && npm run test
  sonarcloud:
    needs: [backend, android, frontend]
    steps:
      - uses: SonarSource/sonarqube-scan-action@v4
\end{minted}
\end{listing}

\section{Nasazení}

Tato sekce popisuje produkční prostředí systému: zvolenou serverovou infrastrukturu, deployment pipeline, monitoring a~způsob distribuce mobilní aplikace.

Pro provoz systému byl zvolen pronajatý VPS server s~operačním systémem Ubuntu. Oproti cloudovým platformám (AWS, Azure) toto řešení nabízí nižší provozní náklady a~plnou kontrolu nad konfigurací, což odpovídá rozsahu projektu i~technickým možnostem cílové organizace. Webová aplikace je veřejně dostupná na adrese \url{https://majetek.zs-ustavni.cz}. Architekturu nasazení znázorňuje obrázek~\ref{fig:system-architecture}.

\subsection{Serverová infrastruktura}

Serverovou infrastrukturou jsou míněny všechny komponenty běžící na VPS, které dohromady umožňují obsluhu HTTPS požadavku od uživatele -- reverse proxy, backendová aplikace, databáze a~doplňková provozní opatření (zálohy, zabezpečení vzdáleného přístupu). Architektura odpovídá běžnému třívrstvému uspořádání webové aplikace: požadavek z~prohlížeče dorazí na reverse proxy, ta obslouží statické soubory frontendu nebo přesměruje API volání na backend, který přistupuje k~datům v~databázi. Konkrétně na serveru běží tyto tři služby:

\begin{itemize}
	\item \textbf{Nginx} -- reverse proxy, která obsluhuje HTTPS požadavky. Statické soubory frontendu jsou obsluhovány přímo z~adresáře \texttt{/var/www/majetek} s~SPA fallback (\texttt{try\_files}). Požadavky na \texttt{/api/} jsou přesměrovány na backend (\texttt{proxy\_pass http://localhost:8080}). SSL certifikát je vydán autoritou Let's Encrypt s~automatickou obnovou.
	\item \textbf{Spring Boot backend} -- běží jako systemd služba na portu 8080 a~komunikuje s~databází prostřednictvím JDBC. Při výpadku ji systemd automaticky restartuje.
	\item \textbf{PostgreSQL} -- databázový server s~automatickým denním zálohováním prostřednictvím cron úlohy. Zálohy jsou uchovávány lokálně po dobu 30 dnů.
\end{itemize}

Přístup k~serveru přes SSH je zabezpečen nestandardním portem s~autentizací výhradně klíčem (heslo zakázáno). Nginx je nakonfigurován jako \uv{default server}, který pro požadavky bez očekávaného hostname (typicky při přístupu přímo přes IP adresu) vrací HTTP kód 404, čímž neúmyslně zabraňuje skenovacím botům v~indexaci aplikace.

\subsection{Monitoring a~provozní dohled}

Vzhledem k~rozsahu projektu (jedna organizace, řádově desítky uživatelských akcí denně) nebyl nasazen plnohodnotný monitorovací nástroj typu Grafana nebo Prometheus. Provozní dohled zajišťují:

\begin{itemize}
	\item \textbf{systemd journal} -- standardní výstup backendu je archivován s~rotací podle data, přístupný příkazem \texttt{journalctl}.
	\item \textbf{Nginx access a~error logy} -- ve~standardním adresáři \texttt{/var/log/nginx/} pro analýzu HTTP požadavků a~chyb.
	\item \textbf{automatický restart} -- backend se po pádu sám obnoví díky systemd nastavení.
	\item \textbf{denní zálohy databáze} -- pravidelná cron úloha vytváří dump databáze a~uchovává jej lokálně po dobu 30 dnů.
\end{itemize}

\subsection{Proces nasazení}

Backend a~frontend jsou nasazovány společně, ale odlišnými kroky -- backend je dlouhoběžný proces, který je nutné zastavit a~znovu spustit, zatímco frontend je sada statických souborů, kterou stačí synchronizovat do adresáře obsluhovaného Nginxem. Deployment workflow \texttt{deploy.yml} (spouštěný manuálně z~rozhraní GitHubu) provede následující kroky:

\begin{enumerate}
	\item Sestaví backend JAR soubor příkazem \texttt{mvn clean package}.
	\item Sestaví frontend příkazem \texttt{npm run build}.
	\item Nahraje JAR na server prostřednictvím SCP a~restartuje backendovou systemd službu.
	\item Synchronizuje frontendové soubory na server prostřednictvím rsync (přepíše obsah \texttt{/var/www/majetek}); restart Nginxu není nutný, protože soubory jsou obslouženy přímo.
\end{enumerate}

V~případě potřeby zásahu mimo workflow lze server obsloužit ručně přes SSH (např.\ pro úpravu konfigurace nebo obnovu ze~zálohy).

\subsection{Distribuce mobilní aplikace}

Mobilní aplikace je distribuována odděleně jako APK soubor, který se na cílové zařízení Zebra TC22 nahrává ručně prostřednictvím USB. Jedná se o~interní nástroj určený pro jedno konkrétní zařízení, publikace na Google Play proto není nutná.


